\section{A uniform mass gap for positive subhomogeneous additions}
\label{sec:cone}

Fix \(N\ge1\), \(m>1\), and a real-even symbol
\(a\in C^{N+2}(\mathbb R^N\setminus\{0\})\) such that
\begin{equation}\label{cone:principal}
a(t\xi)=t^m a(\xi),\qquad
c|\xi|^m\le a(\xi)\le C|\xi|^m
\end{equation}
for \(t>0\) and positive constants \(c,C\). Put
\[
s=\frac m2,\qquad r=\frac mN,\qquad p=2+2r,
\qquad A_a(f)=\int a(\xi)|\widehat f(\xi)|^2\,d\xi.
\]
The Fourier transform is unitary. Define \(M(f)=\|f\|_2^2\) and let
\(M_{c,a}\) be the sharp critical mass determined by
\begin{equation}\label{cone:gn}
\|f\|_p^p\le\frac p2
\left(\frac{M(f)}{M_{c,a}}\right)^r A_a(f),
\qquad f\in H^{m/2}(\mathbb R^N;\mathbb C).
\end{equation}
The reference mass in this section belongs to the principal
homogeneous symbol \(a\).

\subsection{The two admissible classes}

An added dispersion symbol \(b\) is admissible if it is real and even,
belongs to \(C^{N+2}(\mathbb R^N\setminus\{0\})\), and satisfies
\begin{align}
b(\xi)&\ge0,\qquad \xi\cdot\nabla b(\xi)\le m b(\xi),
&&\xi\ne0,\label{cone:dispersion-sign}\\
|\partial^\beta b(\xi)|
&\le C_{b,\beta}|\xi|^{-|\beta|}(1+|\xi|^m),
&&|\beta|\le N+2,\quad \xi\ne0.\label{cone:dispersion-regularity}
\end{align}
Each constant \(C_{b,\beta}\) must be finite for its individual symbol;
there is no uniform bound on these constants across the class.
There is no lower bound on \(\xi\cdot\nabla b\). Values at the single
frequency zero do not affect the multiplier.

An added nonlinear energy density \(G\) is admissible if
\begin{equation}\label{cone:nonlinear-sign}
G\in C^1([0,\infty)),\quad G(0)=0,\quad G\ge0,\qquad
0\le\rho G'(\rho)\le(1+r)G(\rho)\quad(\rho>0),
\end{equation}
and the map
\begin{equation}\label{cone:nonlinear-regularity}
z\longmapsto g(|z|^2)z,\qquad g=2G',
\end{equation}
is locally \(C^1\) on \(\mathbb C\simeq\mathbb R^2\).
This last requirement is separate from the inequalities in
\eqref{cone:nonlinear-sign}; it supplies the Sobolev chain rule.
These assumptions imply, for a finite constant depending on \(G\),
\begin{equation}\label{cone:nonlinear-growth}
0\le G(\rho)\le C_G(\rho+\rho^{1+r}),\qquad
0\le g(\rho)\le C_G(1+\rho^r).
\end{equation}
For \(\rho\le1\) this follows from the continuity of \(G'\) and
\(G(0)=0\). For \(\rho\ge1\), integrate
\(\rho G'(\rho)\le(1+r)G(\rho)\), or apply the differential inequality
directly if \(G(1)=0\), and then use it to bound \(G'\).

Write
\[
B_b(f)=\int b(\xi)|\widehat f(\xi)|^2\,d\xi,\qquad
F_G(f)=\int G(|f|^2)\,dx,
\]
and define
\begin{equation}\label{cone:energy}
E_{a,b,G}(f)=\frac12\bigl(A_a(f)+B_b(f)\bigr)
-\frac1p\|f\|_p^p+F_G(f).
\end{equation}
All these expressions are finite on \(H^{m/2}\) by
\eqref{cone:dispersion-regularity},
\eqref{cone:nonlinear-growth}, and Sobolev interpolation.

\begin{theorem}[Uniform mass gap over the full addition classes]
\label{cone:gap}
There exists \(\varepsilon_{N,m,a}>0\) such that, for every admissible
\(b,G\), every \(\omega\in\mathbb R\), and every \(v\in\mathbb R^N\),
each nonzero weak \(H^{m/2}\) solution of
\begin{equation}\label{cone:profile}
(a(D)+b(D))U+\omega U+iv\cdot\nabla U
-|U|^{2r}U+g(|U|^2)U=0
\end{equation}
with \(E_{a,b,G}(U)<0\) satisfies
\[
M(U)\ge M_{c,a}+\varepsilon_{N,m,a}.
\]
The gap is independent of the coefficients, derivative bounds,
and detailed shapes of the admissible additions. It is also
independent of frequency and velocity.
\end{theorem}

The theorem concerns traveling profiles. It does not assert their
existence for every admissible model, nor does it assert a
decomposition or classification of arbitrary global solutions.
The proof does not require radiality, finite spatial moments,
uniqueness of the principal optimizer, or real-sector
nondegeneracy of that optimizer.

\subsection{Regularity and the cutoff dilation identity}

\begin{lemma}[Individual profile regularity]\label{cone:regularity}
Every weak \(H^{m/2}\) profile in \eqref{cone:profile} belongs to
\(L^\infty\cap H^{m+1}\).
\end{lemma}
\begin{proof}
For the fixed symbol \(b\) and velocity \(v\), choose \(\Lambda>0\)
large enough that
\[
a(\xi)+b(\xi)-v\cdot\xi+\Lambda\asymp1+|\xi|^m.
\]
Here \(b\ge0\), ellipticity, and \(m>1\) give the lower bound,
and \eqref{cone:dispersion-regularity} gives the upper bound.
The resolvent
\[
R_\Lambda=(a(D)+b(D)+iv\cdot\nabla+\Lambda)^{-1}
\]
maps \(L^q\) into the Bessel potential space \(H^{m,q}\) for every
\(1<q<\infty\). Indeed the multiplier
\[
\frac{\langle\xi\rangle^m}
{a(\xi)+b(\xi)-v\cdot\xi+\Lambda}
\]
satisfies the Mihlin bounds by
\eqref{cone:dispersion-regularity} and homogeneity of \(a\).
Constants may depend on the individual model and velocity.

The equation takes the form
\[
U=R_\Lambda\bigl(|U|^{2r}U-g(|U|^2)U\bigr)
+(\Lambda-\omega)R_\Lambda U.
\]
The nonlinear expression has magnitude at most
\(C(|U|+|U|^{1+2r})\) by \eqref{cone:nonlinear-growth}.
For completeness, when \(N>m\), start with
\(q_0=2N/(N-m)\) and reciprocal exponent
\(\beta_0=1/2-m/(2N)\). On \(\{|U|>1\}\) the nonlinear
expression is bounded by \(C|U|^{1+2r}\), and on
\(\{|U|\le1\}\) by \(C|U|\). Apply the resolvent to these
two terms separately. Sobolev embedding and interpolation give
the improvement
\[
\beta_{j+1}=(1+2r)\beta_j-r
\]
until an exponent exceeds the threshold for boundedness. The
linear contribution has the better reciprocal output exponent
\(\beta_j-r\) and remains in \(L^2\); interpolation gives a common
output exponent for the two terms. The ideal iterates are
\(\beta_j=1/2-(r/2)(1+2r)^j\), so finitely many iterations
suffice. All nonlinear input exponents exceed one since
\(q_0>1+2r\), and intermediate exponents can be weakened slightly
to avoid endpoints. If \(N=m\), the initial energy space embeds
in every finite \(L^q\), and one starts with sufficiently large
\(q\); if \(N<m\), it already embeds into \(L^\infty\).
Thus \(U\in L^\infty\). The right-hand nonlinearities are now in
\(L^2\), giving \(U\in H^m\subset H^1\). The local \(C^1\)
assumption \eqref{cone:nonlinear-regularity} and the corresponding
property of \(z\mapsto|z|^{2r}z\) imply that both nonlinearities
belong to \(H^1\). A final resolvent application gives
\(U\in H^{m+1}\).
\end{proof}

Choose \(\chi\in C_c^\infty(\mathbb R^N;\mathbb R)\) equal to one
near zero, and put
\[
f_R(x)=x\chi(x/R)=R f(x/R),\qquad
L_R=f_R\cdot\nabla+\frac12\operatorname{div}f_R,\qquad R\ge1.
\]
The operator \(L_R\) is skew-symmetric on \(L^2\).

\begin{lemma}[Nonhomogeneous cutoff commutator]\label{cone:commutator}
The commutators \([b(D),L_R]\) are uniformly bounded in \(R\)
as forms \(H^{m/2}\to H^{-m/2}\). For every \(u\in H^{m/2}\),
\begin{equation}\label{cone:commutator-limit}
\langle u,[b(D),L_R]u\rangle
\longrightarrow
\int \xi\cdot\nabla b(\xi)|\widehat u(\xi)|^2\,d\xi.
\end{equation}
The bound may depend on the individual symbol \(b\).
\end{lemma}
\begin{proof}
Apart from the fixed Fourier normalization, the kernel is
\[
\frac i2\bigl(b(\xi)-b(\eta)\bigr)
(\xi+\eta)\cdot\widehat f_R(\xi-\eta).
\]
The zeroth and first derivative bounds in
\eqref{cone:dispersion-regularity} imply
\begin{equation}\label{cone:kernel-estimate}
|b(\xi)-b(\eta)|\,|\xi+\eta|
\le C_b|\xi-\eta|\bigl(1+(|\xi|+|\eta|)^m\bigr).
\end{equation}
To see this, write \(\rho=|\xi|+|\eta|\). If
\(|\xi-\eta|\ge\rho/4\), use the zeroth-order bound.
Otherwise the segment between the two frequencies has magnitude
comparable to \(\rho\), and the first derivative bound and
mean-value theorem apply.

For \(k=\xi-\eta\), divide \eqref{cone:kernel-estimate} by
\(\langle\xi\rangle^{m/2}\langle\eta\rangle^{m/2}\).
The result is bounded by \(C_b|k|\langle k\rangle^{m/2}\);
this follows from the inequalities for
\(\langle\xi\rangle/\langle\eta\rangle\) and its reciprocal
by a constant times \(\langle k\rangle\).
Since \(\widehat f_R(k)=R^{N+1}\widehat f(Rk)\),
\[
\int |k|\langle k\rangle^{m/2}|\widehat f_R(k)|\,dk
=\int |\ell|\langle\ell/R\rangle^{m/2}
|\widehat f(\ell)|\,d\ell\le C_f.
\]
Young's inequality proves the uniform form bound.

For Schwartz \(u\), \(L_Ru\to Lu\) in \(L^2\), where
\(L=x\cdot\nabla+N/2\). Thus the left side of
\eqref{cone:commutator-limit} converges to
\(2\operatorname{Re}\langle b(D)u,Lu\rangle\).
Fourier integration by parts gives the asserted integral. One may
first exclude a ball of radius \(\delta\) around zero: \(b\)
is bounded there by \eqref{cone:dispersion-regularity}, so the
boundary term is \(O(\delta^N)\); the term at infinity vanishes
by Schwartz decay. Finally
\[
|\xi\cdot\nabla b(\xi)|\le C_b(1+|\xi|^m)
\]
by the same derivative assumption. Density and the uniform form
bound extend the limit to \(H^{m/2}\). This argument does not use
a lower radial-slope bound.
\end{proof}

The identical argument applies to \(a\), giving limit \(mA_a(u)\)
by homogeneity. Applied to \(iv\cdot\nabla\), it gives limit
\(-v\cdot P(u)\), where
\[
P(u)=\int \xi|\widehat u(\xi)|^2\,d\xi.
\]
This momentum is finite because \(m>1\). The degree-one
commutator is bounded on \(H^{1/2}\), and therefore on the
energy space.

\subsection{Exact defects and the shifted frequency}

For a profile \(U\), abbreviate
\[
M=M(U),\quad A=A_a(U),\quad B=B_b(U),\quad
F=F_G(U),\quad E=E_{a,b,G}(U),\quad \rho=|U|^2.
\]
Define the two dilation defects
\begin{align}
D_b&=\int\bigl(mb(\xi)-\xi\cdot\nabla b(\xi)\bigr)
|\widehat U(\xi)|^2\,d\xi\ge0,\label{cone:defect-b}\\
K_G&=N\int\bigl(\rho G'(\rho)-G(\rho)\bigr)\,dx,\qquad
D_G=mF-K_G\ge0.\label{cone:defect-G}
\end{align}
Each defect is finite. In fact
\(-NF\le K_G\le mF\), so \(0\le D_G\le(m+N)F\).
There is no required estimate of \(D_b\) by \(B\), uniformly
over the dispersion class.

\begin{lemma}[Profile identities]\label{cone:identities}
Every profile obeys
\begin{align}
v\cdot P(U)&=2mE-D_b-2D_G,\label{cone:pohozaev}\\
\omega M&=r(A+B)+(2m-p)E-D_b-\vartheta_ND_G,
\qquad \vartheta_N=2(1-1/N).\label{cone:frequency}
\end{align}
Writing \(U=X+iY\) with real \(X,Y\), one has
\begin{equation}\label{cone:phase-energy}
mE=\left\langle Y,
\bigl[a(D)+b(D)+\omega-|U|^{2r}+g(|U|^2)\bigr]Y\right\rangle
+\frac{D_b}{2}+D_G.
\end{equation}
\end{lemma}
\begin{proof}
Lemma~\ref{cone:regularity} makes \(L_RU\) an admissible test.
The kinetic and drift terms converge by
Lemma~\ref{cone:commutator}. The frequency pairing vanishes.
The critical focusing term converges to
\((m/p)\|U\|_p^p\). For the added nonlinearity, the chain rule gives
\(\nabla G(|U|^2)=g(|U|^2)\operatorname{Re}(\overline U\nabla U)\),
and compact-support integration by parts yields
\[
\operatorname{Re}\int g(\rho)U\,\overline{L_RU}\,dx
=\int\operatorname{div}f_R\bigl(\rho G'(\rho)-G(\rho)\bigr)\,dx
\longrightarrow K_G.
\]
Its integrand is dominated by a constant times \(G(\rho)\).
These limits give
\[
\frac12\left(mA+\int\xi\cdot\nabla b\,|\widehat U|^2\right)
-\frac12v\cdot P(U)-\frac mp\|U\|_p^p+K_G=0,
\]
which is \eqref{cone:pohozaev}.

Pairing the profile with \(U\) gives
\[
A+B+\omega M-v\cdot P(U)-\|U\|_p^p+J=0,
\qquad J=\int g(\rho)\rho\,dx.
\]
Since
\[
J=2F+\frac2N K_G,\qquad
pF-J=\frac2N D_G,\qquad
\|U\|_p^p=\frac p2(A+B)+pF-pE,
\]
substitution proves \eqref{cone:frequency}.

Real-even symbols preserve real functions, so the imaginary
profile equation reads
\[
\bigl[a(D)+b(D)+\omega-|U|^{2r}+g(|U|^2)\bigr]Y
=-v\cdot\nabla X.
\]
Also \(P(U)=-2\int Y\nabla X\). Pairing this equation with \(Y\)
and using \eqref{cone:pohozaev} proves
\eqref{cone:phase-energy}.
\end{proof}

The useful frequency is
\begin{equation}\label{cone:shifted-frequency}
\widetilde\omega=\omega+\frac{D_b+\vartheta_ND_G}{M},
\qquad
\widetilde\omega M=r(A+B)+(2m-p)E.
\end{equation}
Substituting this frequency into \eqref{cone:phase-energy} gives
the exact identity
\begin{equation}\label{cone:shifted-identity}
\begin{aligned}
mE={}&
\left\langle Y,\bigl[a(D)+\widetilde\omega-|U|^{2r}\bigr]Y\right\rangle
+B_b(Y)+\int g(|U|^2)Y^2\,dx\\
&+D_b\left(\frac12-\frac{M(Y)}M\right)
+D_G\left(1-\vartheta_N\frac{M(Y)}M\right).
\end{aligned}
\end{equation}
This identity retains the compensating contributions of the
defects. In particular, an arbitrarily large negative shift of
the actual frequency need not be estimated separately.

\subsection{Proof of the uniform gap}

\begin{proof}[Proof of Theorem~\ref{cone:gap}]
The sharp inequality and positivity of \(B,F\) imply that
negative energy requires \(M>M_{c,a}\). If the asserted gap
failed, choose admissible additions \(b_n,G_n\) and profiles
\(U_n\) such that \(E_n<0\) and \(M(U_n)\to M_{c,a}\).

Critical dilation acts on the entire model by
\begin{align*}
U_\lambda(x)&=\lambda^{N/2}U(\lambda x),&
b_\lambda(\xi)&=\lambda^m b(\xi/\lambda),\\
G_\lambda(\rho)&=\lambda^{m+N}G(\lambda^{-N}\rho),&
g_\lambda(\rho)&=\lambda^m g(\lambda^{-N}\rho),\\
\omega_\lambda&=\lambda^m\omega,&
v_\lambda&=\lambda^{m-1}v.
\end{align*}
It preserves both admissible classes and the profile equation,
preserves mass, and multiplies the energy by \(\lambda^m\).
Thus normalize \(A_a(U_n)=1\). The individual regularity constants
may change under scaling, but no uniform bound on them is used.

Set
\[
\delta_n=\left(\frac{M(U_n)}{M_{c,a}}\right)^r-1\longrightarrow0.
\]
The sharp inequality and negative energy give
\begin{equation}\label{cone:near-threshold}
0\le B_n+2F_n<\delta_n,\qquad
-\frac{\delta_n}{2}\le E_n<0,\qquad
\|U_n\|_p^p\longrightarrow\frac p2.
\end{equation}
Lemma~\ref{foundation:compact} gives translations and constant
phases for which \(U_n\to Q\) strongly in \(H^{m/2}\), where
\(Q\) is a real optimizer normalized by
\[
M(Q)=M_{c,a},\qquad A_a(Q)=1,\qquad
\omega_Q=r/M_{c,a}.
\]
Equation \eqref{cone:shifted-frequency} gives
\(\widetilde\omega_n\to\omega_Q\).
After a further constant phase tending to zero, arrange
\(\langle Y_n,Q\rangle=0\), where \(Y_n=\operatorname{Im}U_n\).

By Lemma~\ref{foundation:phase}, the principal phase form
\(L_-=a(D)+\omega_Q-|Q|^{2r}\) is coercive on real functions
\(L^2\)-orthogonal to \(Q\). Strong \(L^p\) convergence gives
\[
|U_n|^{2r}\longrightarrow|Q|^{2r}
\quad\hbox{in }L^{p/(2r)}.
\]
H\"older's inequality with two \(L^p\) factors and Sobolev
embedding transfer this coercivity to the shifted forms:
\[
\left\langle Y_n,
\bigl[a(D)+\widetilde\omega_n-|U_n|^{2r}\bigr]Y_n\right\rangle
\ge\frac{c_Q}{2}\|Y_n\|_{H^{m/2}}^2
\]
for all sufficiently large \(n\).

Moreover \(M(Y_n)/M(U_n)\to0\). Once this ratio is at most \(1/4\),
and since \(0\le\vartheta_N<2\), identity
\eqref{cone:shifted-identity} gives
\[
\begin{aligned}
mE_n\ge{}&
\frac{c_Q}{2}\|Y_n\|_{H^{m/2}}^2+B_{b_n}(Y_n)
+\int g_n(|U_n|^2)Y_n^2\,dx\\
&+\frac14D_{b,n}+\frac12D_{G,n}\ge0.
\end{aligned}
\]
All displayed terms are finite and nonnegative; for the
multiplication term use
\(\int g_n(|U_n|^2)Y_n^2\le J_n\le2(1+r)F_n\).
This contradicts \(E_n<0\). No boundedness of the actual
frequencies, and no convergence or boundedness of \(D_{b,n}\),
was needed.
\end{proof}

\subsection{Uniformity under changes of the principal symbol}

\begin{theorem}[Local uniform robustness]\label{cone:robustness}
Fix a principal symbol \(a_0\) satisfying
\eqref{cone:principal}. There are \(\eta>0\) and
\(\varepsilon>0\) such that the conclusion
\[
M(U)\ge M_{c,a}+\varepsilon
\]
holds for every individually admissible principal symbol \(a\)
with
\[
\sup_{\xi\in S^{N-1}}\left|\frac{a(\xi)}{a_0(\xi)}-1\right|<\eta,
\]
simultaneously for every admissible pair \(b,G\).
No uniform derivative bounds on any symbol are required.
\end{theorem}
\begin{proof}
If false, choose \(a_n\to a_0\) uniformly relative to \(a_0\)
on the sphere, admissible additions, and negative-energy profiles
whose masses differ from \(M_{c,a_n}\) by a quantity tending to
zero. The principal form comparison
\[
(1-o(1))A_{a_0}\le A_{a_n}\le(1+o(1))A_{a_0}
\]
implies convergence of sharp constants and hence
\(M_{c,a_n}\to M_{c,a_0}\). Normalize \(A_{a_n}(U_n)=1\).
The bounds \eqref{cone:near-threshold} hold with \(M_{c,a_n}\);
in addition \(A_{a_0}(U_n)\to1\). Variational compactness for
\(a_0\) gives convergence to its real normalized optimizer \(Q\).
The shifted frequency again converges to \(r/M_{c,a_0}\).
The principal forms converge in norm from \(H^{m/2}\) to
\(H^{-m/2}\), so their phase coercivity transfers exactly as in
the preceding proof. The added forms and dilation defects are
retained as nonnegative terms in
\eqref{cone:shifted-identity}. The same contradiction follows.
\end{proof}

\begin{proposition}[Compact principal families]\label{cone:compact-family}
Let the restrictions to \(S^{N-1}\) of a family of principal
symbols form a compact subset of \(C^0(S^{N-1})\), with every
member satisfying \eqref{cone:principal}. A single positive gap
works over that family and over both full addition classes.
\end{proposition}
\begin{proof}
Apply Theorem~\ref{cone:robustness} at every member, extract
a finite subcover, and take the minimum of the corresponding
positive gap constants.
\end{proof}

\subsection{Examples and the structural mechanism}

For \(a(\xi)=|\xi|^4\), the theorem includes
\(b(\xi)=\mu|\xi|^2\) for every \(\mu\ge0\), with one common
gap independent of \(\mu\). More generally it includes arbitrary
finite nonnegative sums of real-even nonnegative homogeneous
symbols of orders \(0<k\le m\), and nonnegative constant terms.
The added symbols need not be elliptic.

Nonhomogeneous examples include
\[
b(\xi)=\mu\frac{|\xi|^k}{1+|\xi|^k}\quad(0<k\le m),\qquad
b(\xi)=\mu e^{-|\xi|^2},\qquad
b(\xi)=\mu(1+|\xi|^2)^{-q},
\]
simultaneously for arbitrary \(\mu,q\ge0\). The Gaussian has
logarithmic radial slope \(-2|\xi|^2\), which is unbounded below;
this illustrates why the shifted-frequency argument is stronger
than imposing a two-sided radial-slope estimate.
Nonnegative smooth radially decreasing compactly supported
frequency multipliers are also admissible.

The nonlinear class includes all finite sums
\[
G(\rho)=\sum_j\frac{\kappa_j}{q_j+2}\rho^{1+q_j/2},
\qquad \kappa_j\ge0,\qquad 0\le q_j\le2m/N,
\]
which add \(\sum_j\kappa_j|U|^{q_j}U\) to the profile equation.
It also includes, for arbitrary \(\kappa\ge0\),
\[
G(\rho)=\kappa\frac{\rho^{1+r}}{(1+\rho)^r},\qquad
G(\rho)=\kappa\log(1+\rho),\qquad
G(\rho)=\kappa\frac{\rho}{1+\rho}.
\]
The latter two are concave energy densities. Thus neither
convexity of \(G\) nor monotonicity of \(g=2G'\) is required.

The argument uses two complementary consequences of being near
the sharp principal mass. Negative energy forces the actual added
energies \(B+2F\) to be small after kinetic normalization, even if
their coefficients are arbitrarily large. Variational compactness
then makes the nonconstant phase component small. In the shifted
identity, that small phase mass allows the positive dilation
defects to dominate their own negative contributions to the
frequency. This is a sign-and-scaling mechanism rather than a
small-operator perturbation argument.
